\documentclass[10pt,a4paper]{article}

%% paquetes
\usepackage[utf8]{inputenc}
\usepackage{a4wide,color,verbatim,Sweave,url,xargs,amsmath,booktabs,longtable,eurosym}
\usepackage{tikz}
\usepackage{pgfplots}
\pgfplotsset{compat=1.18}
\usetikzlibrary{arrows,calc,positioning}

%% nuevos entornos
\newenvironment{question}{\item}{}
\newenvironment{solution}{\comment}{\endcomment}
\newenvironment{answerlist}{\renewcommand{\labelenumii}{(\alph{enumii})}\begin{enumerate}}{\end{enumerate}}

%% párrafos
\setlength{\parskip}{0.7ex plus0.1ex minus0.1ex}
\setlength{\parindent}{0em}

\begin{document}
\Sconcordance{concordance:PuntoK_Trigonometria_v1.tex:PuntoK_Trigonometria_v1.Rnw:1 %
23 1 1 250 4 1 1 6 5 1 1 6 2 1 1 6 3 1 1 6 2 1 1 6 2 1 1 4 4 0 1 2 37 1 %
1 20 4 0 1 2 19 1}


\begin{enumerate}


\begin{question}

Un punto K se mueve sobre un segmento QT. El <U+00E1>ngulo $\alpha$ y la medida $h$ se relacionan mediante la raz<U+00F3>n trigonom<U+00E9>trica $sin(alpha) = frac{h}{KP}$, de donde se deduce la distancia entre $S$ y $P$ como $SP = \frac{h}{\sin(\alpha)}$ o $SP = h \times \csc(\alpha)$.

\includegraphics[width=8cm]{diagrama.png}
<U+00BF>Cu<U+00E1>l es la gr<U+00E1>fica que muestra las distancias $SP$, cada vez que $S$ se mueve sobre el segmento $QT$?

\begin{tabular}{cc}
A. & B. \\

\includegraphics[width=4cm]{grafica_A.png}
&

\includegraphics[width=4cm]{grafica_B.png}
\\
C. & D. \\

\includegraphics[width=4cm]{grafica_C.png}
&

\includegraphics[width=4cm]{grafica_D.png}
\end{tabular}

\begin{answerlist}
  \item A
  \item B
  \item C
  \item D
\end{answerlist}
\end{question}

\begin{solution}

La relaci<U+00F3>n trigonom<U+00E9>trica dada es $sin(alpha) = frac{h}{KP}$, lo que implica que $SP = \frac{h}{\sin(\alpha)}$.

\textbf{An<U+00E1>lisis matem<U+00E1>tico del comportamiento de la funci<U+00F3>n:}

\begin{itemize}
\item Cuando $\alpha$ se acerca a $0<U+00B0>$: $\sin(\alpha) \to 0$, por lo tanto $SP = \frac{h}{\sin(\alpha)} \to +\infty$ (la distancia tiende a infinito).

\item Cuando $\alpha = 30<U+00B0>$: $\sin(30<U+00B0>) = 0.5$, por lo tanto $SP = \frac{h}{0.5} = 2h$.

\item Cuando $\alpha = 90<U+00B0>$: $\sin(90<U+00B0>) = 1$, por lo tanto $SP = \frac{h}{1} = h$ (valor m<U+00ED>nimo posible).

\item La funci<U+00F3>n $f(\alpha) = \frac{h}{\sin(\alpha)}$ es estrictamente \textbf{decreciente} en el intervalo $(0<U+00B0>, 90<U+00B0>)$ porque:
  \begin{itemize}
  \item $\sin(\alpha)$ es creciente en $(0<U+00B0>, 90<U+00B0>)$
  \item $\frac{1}{\sin(\alpha)}$ es decreciente cuando $\sin(\alpha)$ es creciente
  \item Por lo tanto, $h \cdot \frac{1}{\sin(\alpha)}$ es decreciente
  \end{itemize}
\end{itemize}

\textbf{An<U+00E1>lisis de las opciones:}

\begin{itemize}
\item \textbf{Opci<U+00F3>n A}: Funci<U+00F3>n constante - \textbf{Incorrecta}. La distancia $SP$ no permanece constante cuando $\alpha$ var<U+00ED>a.

\item \textbf{Opci<U+00F3>n B}: Funci<U+00F3>n decreciente - \textbf{Correcta}. Coincide con el an<U+00E1>lisis matem<U+00E1>tico de $\frac{h}{\sin(\alpha)}$.

\item \textbf{Opci<U+00F3>n C}: Funci<U+00F3>n que crece y luego decrece - \textbf{Incorrecta}. No hay raz<U+00F3>n matem<U+00E1>tica para que la funci<U+00F3>n tenga un m<U+00E1>ximo local.

\item \textbf{Opci<U+00F3>n D}: Funci<U+00F3>n constante - \textbf{Incorrecta}. Similar al error de la opci<U+00F3>n A.
\end{itemize}

La respuesta correcta es: \textbf{D}.

\begin{answerlist}
  \item Falso. Esta opci<U+00F3>n no representa correctamente el comportamiento de la funci<U+00F3>n KP = h/sin(<U+03B1>).
  \item Falso. Una funci<U+00F3>n constante implicar<U+00ED>a que KP no depende de <U+03B1>, lo cual contradice la relaci<U+00F3>n trigonom<U+00E9>trica dada.
  \item Falso. La funci<U+00F3>n h/sin(<U+03B1>) no tiene m<U+00E1>ximos locales en el intervalo (0<U+00B0>, 90<U+00B0>), es mon<U+00F3>tona decreciente.
  \item Verdadero. La funci<U+00F3>n KP = h/sin(<U+03B1>) es estrictamente decreciente en (0<U+00B0>, 90<U+00B0>).
\end{answerlist}
\end{solution}

% Meta-information ICFES
\exname{Trigonometr<U+00ED>a - Interpretaci<U+00F3>n de Gr<U+00E1>ficas}
\extype{schoice}
\exsolution{0001}
\exshuffle{TRUE}
\exsection{Geometr<U+00ED>a - Trigonometr<U+00ED>a}

% Metadatos ICFES adicionales (como comentarios)
% Competencia: interpretacion_representacion
% Nivel de dificultad: 3
% Componente: geometrico_metrico
% Contexto: matematico
% Categor<U+00ED>a: geometria
% Tipo: generico

\end{enumerate}
\end{document}
